\chapter{Як машина вчиться рухатися}
\chaptermark{Як машина вчиться рухатися}
\label{sec:motorlearning}
\begin{chaptergoals}
\item чому рухам потрібен третій учитель --- провід помилки до кожного виходу --- і чого він коштує;
\item як \nt{GLUT}-розширювач і провід помилки на кожен вихід дають правило найменших квадратів;
\item як навчене коло стає внутрішньою моделлю --- адаптивним фільтром, що передбачає тіло;
\item чому швидкий і повільний процеси навчання пояснюють післядію і спонтанне відновлення.
\end{chaptergoals}

\section{Три вчителі}

У розділі~\ref{sec:muscles} ми бачили, як машина рухає м'язами, і зупинилися на гіпотезі: швидкі й точні рухи потребують внутрішньої моделі тіла, яка передбачає результат команди~\cite{WolpertGhahramani2000}. Звідки береться ця модель? Її треба вивчити, і вивчити не так, як ми вчилися досі.

У книзі вже було два вчителі. Перший --- \emph{жодного}: правило Гебба (розділ~\ref{sec:dofa}) посилює те, що часто буває разом, і стискає входи. Другий --- \emph{винагорода}: \nt{DOPA} розсилає всім одне число, «краще чи гірше, ніж очікувалося». Рухові потрібен третій учитель --- \emph{помилка}. Рука промахнулася повз чашку на два сантиметри ліворуч: це не «гірше», це вектор, який каже кожному виходу, куди і наскільки він помилився. Інженер назвав би це навчанням з учителем.

Різниця між другим і третім учителем --- це різниця між одним проводом на всіх і окремим проводом на кожен вихід. У твердженні~\ref{prop:broadcastcost} навчання за одним спільним числом сповільнювалося з кожним нейроном, чий шум впливав на результат. Якщо ж кожен вихід $i$ отримує свою помилку $e_i$, ваги можна змінювати за найпростішим правилом:
\begin{empheq}[box=\keybox]{equation}
\frac{\dd w_{ij}}{\dd t} \;=\; -\eta\,e_i(t)\,x_j(t) .
\label{eq:lms}
\end{empheq}
Це правило найменших квадратів, давно відоме в техніці адаптивних фільтрів~\cite{WidrowHoff1960}: вага змінюється пропорційно добуткові свого входу на помилку свого виходу й у середньому йде за градієнтом квадрата помилки. Жодного шуму-пошуку, як у розділі~\ref{sec:dofa}, не потрібно: напрямок поправки надходить просто проводом помилки. І швидкість не падає з кількістю виходів, бо чужі помилки до синапсу не доходять.

\begin{keyblock}
\begin{proposition}[Три вчителі]
\label{prop:teachers}
Правило Гебба без учителя навчає того, що часте; сигнал \nt{DOPA} одним числом на всю мережу навчає того, що вигідне, і тим повільніше, що більше нейронів; провід помилки до кожного виходу навчає за правилом~\eqref{eq:lms} того, що точне, і швидкість не залежить від кількості виходів. Ціна третього способу --- окремий провід на кожен вихід і хтось, хто знає правильну відповідь.
\end{proposition}
\end{keyblock}

\section{Провід помилки}

Хто може знати правильну відповідь для руху? Самі датчики. Якщо рука пішла ліворуч від цілі, про це повідомляють око й датчики розтягу з розділу~\ref{sec:sensors}. Потрібна схема, яка доставить цю помилку до потрібних виходів.

У мозку є ділянка, яка, схоже, влаштована саме під це~\cite{Marr1969,Albus1971}. Її схема незвично правильна, майже кристалічна, і в ній легко впізнати правило~\eqref{eq:lms} (рис.~\ref{fig:errorwire}).

\emph{Розширювач входу.} Сигнали про команду й про стан тіла надходять на величезний шар маленьких \nt{GLUT}-нейронів. Їх близько сімдесяти мільярдів --- більше, ніж у всьому іншому мозку разом узятому~\cite{Azevedo2009}. Кожен із них змішує кілька входів, і сигнал розкладається в дуже довгий вектор. Інженерові цей прийом знайомий: якщо підняти розмірність входу, у новому просторі майже будь-яку потрібну залежність можна отримати простою зваженою сумою~\cite{Cover1965}.

\emph{Вихідні нейрони.} Зважену суму обчислюють близько тридцяти мільйонів великих вихідних нейронів~\cite{Andersen1992cereb}, у кожного --- порядку ста тисяч входів від розширювача. І це \nt{GABA}-нейрони: результат навчання передається далі не збудженням, а гальмуванням, як у забороні розділу~\ref{sec:gaba}.

\emph{Провід помилки.} Кожен вихідний нейрон отримує ще один вхід, особливий: рівно один \nt{GLUT}-провід, що спрацьовує рідко, близько разу на секунду, але щоразу викликає у вихідному нейроні потужний спалах~\cite{Ito2001}. Це і є $e_i$: імовірність спалаху залежить від напрямку помилки руху~\cite{Herzfeld2018}. Збіг спалаху помилки з активністю входу послаблює цей вхід --- точнісінько член $-\eta\,e_i x_j$ у~\eqref{eq:lms}.

\begin{figure}[t]
\centering
\begin{tikzpicture}[>=Stealth,
  gran/.style={circle,draw,thick,fill=blue!6,minimum size=0.32cm,inner sep=0pt},
  outn/.style={circle,draw,very thick,fill=red!10,minimum size=1.1cm,inner sep=0pt,font=\scriptsize},
  inh/.style={-{Bar[width=7pt]},very thick,red!70!black}]
% inputs
\foreach \y/\lab in {1.6/{команда},0.4/{стан тіла}}{
  \draw[->,thick,blue!60!black] (-0.6,\y) node[left,font=\scriptsize,black] {\lab} -- (0.35,\y);}
% expansion layer
\foreach \i in {0,...,11}{\node[gran] (g\i) at ({0.8+0.28*mod(\i,2)},{-0.3+0.23*\i}) {};}
\foreach \i in {0,...,11}{\pgfmathsetmacro{\yy}{mod(\i,3)>0 ? 1.6 : 0.4}\draw[gray!60!black] (0.35,\yy) -- (g\i);}
\node[font=\scriptsize,align=center] at (0.95,-1.05) {розширювач\\$\sim7\cdot10^{10}$ \nt{GLUT}};
% parallel wires to the output neuron
\foreach \i in {0,...,11}{\draw[->,gray!70!black] (g\i.east) -- (4.1,{1.0+0.07*(\i-5.5)});}
\node[outn] (P) at (4.6,1.0) {\nt{GABA}};
\node[font=\scriptsize,align=center,above=2pt] at (P.north) {вихідний нейрон\\$\sim10^5$ входів $x_j$, ваги $w_{ij}$};
\draw[inh] (P) -- (6.8,1.0) node[right,font=\scriptsize,align=left] {поправка\\до руху};
% error wire
\draw[->,very thick,red!70!black,densely dashed] (4.6,-1.4) node[below,font=\scriptsize,black,align=center] {один провід помилки $e_i$\\\nt{GLUT}, $\sim1$ раз/с} -- (P.south);
\end{tikzpicture}
\caption{Схема навчання з учителем, як її, схоже, реалізує мозок. Команда й стан тіла розкладаються величезним шаром \nt{GLUT}-нейронів у довгий вектор $x_j$; вихідний \nt{GABA}-нейрон додає його з вагами $w_{ij}$. Єдиний провід помилки приносить $e_i$; збіг помилки з активним входом послаблює вагу цього входу, як у правилі~\eqref{eq:lms}.}
\label{fig:errorwire}
\end{figure}

Одна трудність знайома нам із розділу~\ref{sec:dofa}: помилка надходить пізніше, ніж вхід, що до неї призвів. Промах видно через десятки--сотні мілісекунд після команди. Отже, синапс має пам'ятати, що був активний, --- потрібен слід, як у правилі~\eqref{eq:threefactor}. І довжина цього сліду, судячи з дослідів, підлаштована під затримку зворотного зв'язку в кожній задачі: там, де помилка надходить через сто мілісекунд, найсильніше послаблюється вхід, що був активним за сто мілісекунд до неї~\cite{Suvrathan2016}.

\begin{keyblock}
\begin{proposition}[Провід помилки]
\label{prop:errorwire}
У схемі, зображеній на рис.~\ref{fig:errorwire}, кожен вихідний нейрон отримує рівно один провід помилки, і збіг помилки з входом послаблює цей вхід. Це правило найменших квадратів~\eqref{eq:lms}, застосоване до входу, розширеного до величезної розмірності. Будову схеми й послаблення за збігу виміряно; що вся схема працює як навчання з учителем, --- провідна гіпотеза.
\end{proposition}
\end{keyblock}

\section{Внутрішня модель як адаптивний фільтр}

Чого вчиться така схема? Найприродніша відповідь --- передбачення. Нехай на вхід надходить команда $u(t)$, пропущена через набір фільтрів із різними сталими часу, як у розширювачі: $x_j(t)=(g_j*u)(t)$. Вихід --- їхня зважена сума, передбачення того, що повідомлять датчики:
\begin{empheq}[box=\keybox]{equation}
\hat y(t) \;=\; \sum_j w_j\,(g_j*u)(t), \qquad e(t) \;=\; \hat y(t) - y(t) ,
\label{eq:adaptivefilter}
\end{empheq}
а ваги вчаться за~\eqref{eq:lms}. Це \emph{адаптивний фільтр}: схема, яка сама добирає свою передавальну функцію так, щоб повторити невідому систему~\cite{Fujita1982}. Тут невідома система --- власне тіло з його масою, пружністю й затримками. Вивчений фільтр і є внутрішньою моделлю розділу~\ref{sec:muscles}: він каже, що повідомлять датчики, не чекаючи на них.

Така модель корисна двічі. По-перше, її передбачення можна подати замість запізнілого зворотного зв'язку, і тоді умова стійкості~\eqref{eq:delaystable} перестає зв'язувати руки: виправляти можна швидко, за моделлю. По-друге, різниця між передбаченим і отриманим --- це сигнал про \emph{несподіване}. Усе, що машина зробила сама, модель передбачила й відняла; лишається те, що зробив світ. Тому, за однією з гіпотез, ми й не можемо полоскотати самі себе~\cite{Blakemore1998}.

\section{Коли світ змінюється: швидке й повільне}

Перевірмо схему дослідом, який легко поставити. Людина тягнеться до цілі, а світ несподівано змінюється: наприклад, на руку діє бічна сила. Перші рухи промахуються, потім промах зменшується. Якщо силу прибрати, рука спершу промахується в \emph{інший} бік --- модель вивчила силу й далі її компенсує. Це пряме свідчення того, що вивчено модель, а не просто «почало виходити».

Досліди показують і тоншу річ: навчаються одночасно два процеси --- швидкий, що швидко вчиться і швидко забуває, і повільний, що вчиться повільно й пам'ятає довго~\cite{Smith2006}. Поправки оновлюються після кожного руху, за подією:
\begin{empheq}[box=\keybox]{equation}
e = p - (x_f + x_s), \qquad x_f \leftarrow A_f\,x_f + B_f\,e, \qquad x_s \leftarrow A_s\,x_s + B_s\,e ,
\label{eq:tworate}
\end{empheq}
де $p$ --- збурення, $x_f$ і $x_s$ --- вивчені поправки двох процесів, $A$ --- наскільки поправка зберігається до наступного руху, $B$ --- наскільки сильно вона вчиться на помилці. У швидкого процесу $B$ велике, а $A$ помітно менше за одиницю; у повільного --- навпаки.

\begin{figure}[t]
\centering
\begin{tikzpicture}[>=Stealth,x=0.034cm,y=1.9cm]
\draw[->,gray!70!black] (0,0) -- (355,0) node[right,font=\small,black] {рух};
\draw[->,gray!70!black] (0,-1.15) -- (0,1.25);
\foreach \v in {-1,1}{\draw[gray!60!black] (-3,\v) -- (3,\v); \node[left,font=\scriptsize] at (0,\v) {$\v$};}
\foreach \n in {100,200,300}{\draw[gray!60!black] (\n,-0.04) -- (\n,0.04); \node[below,font=\scriptsize] at (\n,0) {\n};}
\draw[thin,gray!70!black] plot file {figures/tworate_p.dat};
\draw[thick,orange!85!black] plot file {figures/tworate_fast.dat};
\draw[thick,blue!60!black] plot file {figures/tworate_slow.dat};
\draw[very thick,black] plot file {figures/tworate_net.dat};
\node[font=\scriptsize,gray!50!black] at (120,1.12) {збурення $p$};
\node[font=\scriptsize,black,anchor=west] at (238,0.52) {сума: повернення};
\node[font=\scriptsize,blue!60!black,anchor=west] at (150,0.39) {повільний $x_s$};
\node[font=\scriptsize,orange!85!black,anchor=west] at (60,0.08) {швидкий $x_f$};
\draw[densely dashed,gray!60!black] (226,-1.15) -- (226,1.25);
\node[font=\scriptsize,align=center,anchor=south west] at (228,-1.1) {помилок\\більше немає};
\end{tikzpicture}
\caption{Два процеси навчання за моделлю~\eqref{eq:tworate}, $A_f=0.92$, $B_f=0.1$, $A_s=0.996$, $B_s=0.02$. Двісті рухів зі збуренням $p=1$, далі шість зі зворотним, $p=-1$, доки сума не повернеться до нуля. Після штрихової лінії помилки перестають надходити, і сума сама повертається до старої поправки: швидкий процес забув зворотне збурення, повільний пам'ятає пряме.}
\label{fig:tworate}
\end{figure}

Модель~\eqref{eq:tworate} пояснює дослід, який важко зрозуміти без неї (рис.~\ref{fig:tworate}). У нашому розрахунку за двісті рухів зі збуренням сума поправок сягає $0.84$, причому більшу частину тримає повільний процес, $0.63$. Далі шість рухів зі зворотним збуренням повертають суму до нуля: швидкий процес пішов у мінус, $-0.53$, повільний ще в плюсі, $0.46$. Якщо тепер перестати повідомляти помилки, швидкий процес забуває свою поправку за десятки рухів, повільний --- набагато довше, і сума \emph{сама} піднімається до $0.37$ за сорок рухів: стара навичка повертається без жодного тренування. Таке спонтанне відновлення виміряно в людей; модель з одним процесом його дати не може --- без помилок вона просто згасає до нуля.

Навіщо два процеси? Правдоподібну інженерну відповідь дає оптимальний фільтр розділу~\ref{sec:acet}. Тіло змінюється з різних причин із різною швидкістю: м'язи втомлюються за хвилини, ростуть і слабшають за місяці. Якщо завади бувають швидкими й повільними, найкраща оцінка складається з двох фільтрів із різними сталими часу, і кожен ловить свою~\cite{Kording2007}. Швидкий процес --- це тимчасова поправка, «рука втомилася»; повільний --- «рука стала іншою». Це поки гіпотеза, але вона пояснює, чому навичка, вивчена за один день, частково втрачається до наступного і чому вдруге вона вчиться швидше.

\section{Підсумок}

\begin{table}[t]
\centering
\footnotesize
\setlength{\tabcolsep}{4pt}
\renewcommand{\arraystretch}{1.15}
\begin{tabular}{>{\raggedright}p{3.0cm}>{\raggedright}p{4.6cm}>{\raggedright\arraybackslash}p{5.2cm}}
\toprule
Що робить схема & Як & Що відомо \\
\midrule
вчиться з учителем & провід помилки до кожного виходу, правило~\eqref{eq:lms} & будову схеми й послаблення за збігу виміряно; навчання з учителем --- провідна гіпотеза \\
розширює вхід & сімдесят мільярдів \nt{GLUT}-нейронів & кількість нейронів виміряно; роль розширення --- теорія \\
враховує затримку помилки & слід, підлаштований під затримку & виміряно \\
будує внутрішню модель & адаптивний фільтр~\eqref{eq:adaptivefilter} & добре обґрунтована гіпотеза \\
вчиться двома темпами & швидкий і повільний процеси~\eqref{eq:tworate} & післядію й спонтанне повернення виміряно; два процеси --- модель \\
\bottomrule
\end{tabular}
\caption{Як машина вчиться рухатися.}
\label{tab:motorlearning}
\end{table}

Тепер у машини три вчителі (таблиця~\ref{tab:motorlearning}). Гебб стискає те, що часте; \nt{DOPA} одним числом навчає вибирати вигідне; провід помилки навчає точності, і навчає швидко, бо до кожного виходу йде своя помилка. Мозок, схоже, користується всіма трьома, і кожним там, де він найдешевший: широкомовним сигналом --- для небагатьох рішень, окремими проводами --- для точного підлаштування тіла, де правильну відповідь повідомляють самі датчики.

\section{Задачі}

\begin{exercise}[\lvlA]
\label{ex:15:1}
\emph{Ємність схеми з проводом помилки.} У схемі $3\cdot10^7$ вихідних нейронів по $10^5$ входів, у кожного свій провід помилки ($1$~імп/с). Нейрон з $n$ входами вивчає будь-яку розмітку $\approx2n$ векторів~\cite{Cover1965}. (а)~Скільки асоціацій зберігає схема? (б)~Оцініть за~\eqref{eq:spikeentropy} при $\Delta t=10\ms$ найменший час заповнення ємності нейрона.
\end{exercise}

\begin{exercise}[\lvlB]
\label{ex:15:2}
\emph{Швидкість правила найменших квадратів.} Моди правила~\eqref{eq:lms} згасають зі швидкостями $\eta\lambda_k(C)$. Для п'яти фільтрів~\eqref{eq:adaptivefilter} на білому шумі $C_{jk}=2\sqrt{\tau_j\tau_k}/(\tau_j+\tau_k)$. Знайдіть відношення найшвидшої й найповільнішої швидкостей, якщо $\tau_j$ зростають у $2$ рази ($10$--$160\ms$) і в $4$ рази.
\end{exercise}

\begin{exercise}[\lvlB]
\label{ex:15:3}
\emph{Аналіз двох процесів.} Запишіть модель~\eqref{eq:tworate} з параметрами рис.~\ref{fig:tworate} як $\mathbf x_{n+1}=M\mathbf x_n+\mathbf b\,p$. (а)~Знайдіть за власними числами $M$ сталі часу в рухах. (б)~Знайдіть усталені $x_f$, $x_s$ та їхню суму за сталого $p=1$.
\end{exercise}

\begin{exercise}[\lvlB]
\label{ex:15:4}
\emph{Моделювання: вдруге швидше.} За моделлю~\eqref{eq:tworate} з параметрами з \texttt{models/\allowbreak motor\_learning.py} знайдіть кількість рухів до $x_f+x_s\ge0.8$ при $p=1$: (а)~у ненавченої машини; (б)~після $200$ рухів із $p=1$ і зворотного збурення до нуля суми, як на рис.~\ref{fig:tworate}. Чи може так прискоритися модель з одним процесом?
\end{exercise}

\begin{exercise}[\lvlB]
\label{ex:15:5}
\emph{Регулятор із внутрішньою моделлю.} Об'єкт $\dd x/\dd t=ku$ видно із затримкою $d$; регулятор $u=-G\hat x$ передбачає $\hat x(t)=x(t-d)+\hat k\int_{t-d}^{t}u\,\dd s$. (а)~При $\hat k=k=1$, $d=150\ms$ знайдіть $G$ для сталої часу $20\ms$ і порівняйте з межею~\eqref{eq:delaystable}. (б)~Чи стійка схема при $\hat k=0.4k$?
\end{exercise}

\begin{exercise}[\lvlB]
\label{ex:15:6}
\emph{Адаптивний фільтр навчає м'яз.} Об'єкт --- посмикування~\eqref{eq:twitch} одиничної площі, $T_c=50\ms$; фільтри~\eqref{eq:adaptivefilter} --- RC-кола з $\tau_j=12.5$, $25$, $50$, $100$, $200\ms$; вхід --- нове нормальне число кожні $10\ms$. За скільки секунд правило~\eqref{eq:lms} при $\eta=50$ і $200$ знижує помилку до $0.5\%$? Чому ваги й через $300$~с далекі від найкращих?
\end{exercise}

\begin{exercise}[\lvlC]
\label{ex:15:7}
\emph{Два процеси як оптимальний фільтр.} Збурення --- сума процесів $p^{f,s}_{n+1}=A_{f,s}\,p^{f,s}_n+w^{f,s}_n$ із дисперсіями шумів $q_f$, $q_s$, помилку спостерігають із шумом дисперсії $R$; фільтр Калмана дає~\eqref{eq:tworate}. За яких $q_f/R$, $q_s/R$ з $A_f=0.92$, $A_s=0.996$ виходить $B_f=0.1$, $B_s=0.02$? Як зміняться $B_f$, $B_s$, якщо вчетверо збільшити $q_s$?
\end{exercise}
